%% Copernicus Publications Manuscript Preparation Template for LaTeX Submissions
%% ---------------------------------
%% This template should be used for copernicus.cls
%% The class file and some style files are bundled in the Copernicus Latex Package, which can be downloaded from the different journal webpages.
%% For further assistance please contact Copernicus Publications at: production@copernicus.org
%% https://publications.copernicus.org/for_authors/manuscript_preparation.html


%% Please use the following documentclass and journal abbreviations for preprints and final revised papers.

%% 2-column papers and preprints
\documentclass[journal abbreviation, manuscript]{copernicus}



%% Journal abbreviations (please use the same for preprints and final revised papers)


% Advances in Geosciences (adgeo)
% Advances in Radio Science (ars)
% Advances in Science and Research (asr)
% Advances in Statistical Climatology, Meteorology and Oceanography (ascmo)
% Aerosol Research (ar)
% Annales Geophysicae (angeo)
% Archives Animal Breeding (aab)
% Atmospheric Chemistry and Physics (acp)
% Atmospheric Measurement Techniques (amt)
% Biogeosciences (bg)
% Climate of the Past (cp)
% DEUQUA Special Publications (deuquasp)
% Earth Surface Dynamics (esurf)
% Earth System Dynamics (esd)
% Earth System Science Data (essd)
% E&G Quaternary Science Journal (egqsj)
% EGUsphere (egusphere) | This is only for EGUsphere preprints submitted without relation to an EGU journal.
% Earth Observation (eo)
% European Journal of Mineralogy (ejm)
% Geochronology (gchron)
% Geographica Helvetica (gh)
% Geoscience Communication (gc)
% Geoscientific Instrumentation, Methods and Data Systems (gi)
% Geoscientific Model Development (gmd)
% History of Geo- and Space Sciences (hgss)
% Hydrology and Earth System Sciences (hess)
% Journal of Bone and Joint Infection (jbji)
% Journal of Environmentally Compatible Air Transport System (jecats)
% Journal of Micropalaeontology (jm)
% Journal of Sensors and Sensor Systems (jsss)
% Magnetic Resonance (mr)
% Mechanical Sciences (ms)
% Natural Hazards and Earth System Sciences (nhess)
% Nonlinear Processes in Geophysics (npg)
% Ocean Science (os)
% Polarforschung - Journal of the German Society for Polar Research (polf)
% Proceedings of the International Association of Hydrological Sciences (piahs)
% Proceedings of the International Ocean Drilling Programme (piodp)
% Safety of Nuclear Waste Disposal (sand)
% Scientific Drilling (sd)
% SOIL (soil)
% Solid Earth (se)
% State of the Planet (sp)
% The Cryosphere (tc)
% Weather and Climate Dynamics (wcd)
% Web Ecology (we)
% Wind Energy Science (wes)


%% \usepackage commands included in the copernicus.cls:
\usepackage[german, english]{babel}
\usepackage{tabularx}
\usepackage{cancel}
\usepackage{multirow}
\usepackage{supertabular}
\usepackage{algorithmic}
\usepackage{algorithm}
\usepackage{amsthm}
\usepackage{float}
\usepackage{subfig}
\usepackage{rotating}


\begin{document}

\title{Automated Camera Path Generation for Atmospheric River Visualization in Met3D}


% \Author[affil]{given_name}{surname}

\Author[][]{Samarth Pathak}{samarth.pathak@informatik.hs-fulda.de} %% correspondence author
\Author[][]{Arpana Sitoula}{arpana.sitoula@informatik.hs-fulda.de}
\Author[]{Suraj Kumar Ajampodi}{suraj-kumar.ajampodi@informatik.hs-fulda.de}

\affil[]{Hochschule Fulda,
University of Applied Sciences,
Fachbereich Angewandte Informatik,
Leipziger Straße 123,
36037 Fulda}

%% The [] brackets identify the author with the corresponding affiliation. 1, 2, 3, etc. should be inserted.

%% If an author is deceased, please add \deceased[$Deceased date if applicable$]{$Author number$} (e.g. \deceased[13 November 2015]{2}) at the end of the affiliations. The author number depends on the placement of the author in the author list, e.g. the third author has number 3.


%% If authors contributed equally, please add \equalcontrib{$Author numbers$} (e.g. \equalcontrib{1,3}) at the end of the affiliations. The author number depends on the placement of the author in the author list, e.g. the third author has number 3.




\runningtitle{TEXT}

\runningauthor{TEXT}





\received{}
\pubdiscuss{} %% only important for two-stage journals
\revised{}
\accepted{}
\published{}

%% These dates will be inserted by Copernicus Publications during the typesetting process.


\firstpage{1}

\maketitle

\begin{abstract}
This work presents an automated approach for generating camera paths for tropical
cyclones (TCs) and atmospheric rivers (ARs) in Met3D. Given detected and tracked
weather features from an existing meteorological analysis pipeline, the approach
automatically selects a camera motion for each feature and generates the
corresponding camera path, removing the need to define camera movement manually.

Each tracked feature is characterized by its travel distance, spatial extent, and
shape volatility, and the dominant of these three characteristics determines which
of three camera-motion types is assigned: \textit{follow}, which tracks a feature's
movement; \textit{skeleton}, which traces the spatial structure of an elongated
feature; and \textit{orbit}, which reveals changes in a feature's shape. The
assigned motion type and the feature's tracked positions are then used to generate
a corresponding camera path and its temporal sequence, synchronized with the
advancing forecast data.

The resulting system extends the existing meteorological analysis pipeline with an
automated camera-animation stage: it consumes detected and tracked features and
produces Met3D-compatible camera sequences for the highest-priority features. The
generated sequences are evaluated in Met3D and compared against manually created
sequences with respect to camera movement and generation time.

\textbf{Keywords:} Atmospheric Rivers, Tropical Cyclones, Camera Path Generation,
Automated Camera Planning, Meteorological Visualization, Met3D
\end{abstract}

\section{Introduction}

Met3D is a scientific visualization tool for three-dimensional
meteorological data. This work builds on a pipeline developed in earlier
stages of the project. Live global forecast data from the German Weather
Service is first processed by a detection network. This network labels
each grid cell of the forecast domain as background, tropical cyclone
(TC), or atmospheric river (AR). A separate tracking stage then links
these detections into continuous feature trajectories across the forecast
period. Once detected and tracked, features are displayed in Met3D as an
overlay on the underlying data.

The existing pipeline stops at display. It does not decide what a viewer
should look at, or when. A single global forecast can contain several
active cyclones across different ocean basins, together with atmospheric
rivers spanning an entire hemisphere. A user who wants to inspect all of
these features must currently navigate to each one by hand, stepping
through the forecast period manually. No part of the system identifies
which features deserve attention, and none decides how the camera should
move to present a given feature clearly.

This work adds an automated camera-animation stage to the existing
pipeline. Detected and tracked features are analysed to decide which are
most relevant for visualization, and how each should be filmed. Each
track is characterized by three properties: travel distance, spatial
extent, and shape volatility. Whichever property dominates a given track
determines two things at once: its priority for visualization, and which
of three camera-motion types it is assigned. These three types are
follow, skeleton, and orbit. For elongated features such as atmospheric
rivers, a further step extracts a smooth centreline along the feature's
long axis to guide the skeleton motion. For compact or highly volatile
features, an orbiting motion instead provides several viewing angles onto
the changing structure. The resulting motion assignments and tracked
positions are used to generate a full camera sequence, synchronized with
the advancing forecast time, that Met3D can execute directly.

This work makes the following contributions:
\begin{itemize}
    \item A characterization and ranking scheme that assigns each
    tracked TC or AR feature one of three camera-motion types. The
    assignment is based on which of three quantitative descriptors,
    travel distance, spatial extent, or shape volatility, dominates the
    feature's behaviour, rather than on a single blended score.
    \item Three geometry-aware camera-motion strategies, follow,
    skeleton, and orbit, together with a centreline-extraction method
    for elongated atmospheric-river features.
    \item A mechanism for synchronizing camera motion with the
    advancement of forecast time in Met3D, so that the displayed data
    always matches what the camera is currently observing.
    \item An evaluation of the generated camera sequences against
    manually authored reference sequences.
\end{itemize}

The remainder of this document is organized as follows. Section 2
reviews related work on cyclone and atmospheric-river detection and
tracking, and on camera-path planning for narrative visualization.
Section 3 gives a system-level overview of the pipeline. Section 4
describes how the design evolved through earlier, unsuccessful
approaches. Section 5 describes the methodology, covering feature
characterization and ranking and camera-path generation. Section 6
presents the evaluation, followed by discussion, limitations, and
conclusions.

\section{Related Work}
Closely related to the camera-motion strategies used in this work is a line of
research on automated camera path generation for in-situ visualization of CFD
simulations. Iwata et al.~\cite{iwata2023entropy} propose selecting camera
viewpoints by maximizing information entropy over a rendered image, then
connecting the selected viewpoints with a smooth interpolated path so that the
resulting video remains watchable rather than jumping between viewpoints.
Iwata et al.~\citep{iwata2025smart} extend this approach with quaternion-based
path interpolation between consecutive selected viewpoints, producing a
keyframe sequence of information-rich renderings for later interactive
analysis. Both works select viewpoints using a generic, domain-agnostic
image-quality metric (information entropy of the rendered pixels) that applies
to any scalar or vector field. The present work instead derives camera
placement directly from the semantics of the detected meteorological features
themselves (a cyclone's centroid, an atmospheric river's skeleton), trading
generality for a stronger link between the camera's behaviour and the
physical phenomenon being shown.

A further relevant contribution to camera control is the Toric space of Lino
and Christie \citep{lino2015toric}, a compact representation that expresses
visual composition properties (target size, vantage angle, on-screen position)
directly as camera parameters, enabling efficient automated viewpoint
computation and a viewpoint interpolation technique that preserves these
visual properties continuously along a generated path.

\section{System Overview}

The system consists of four stages arranged as a linear pipeline, with a hand-off from meteorological analysis to cinematographic rendering.

\begin{enumerate}
    \item \textbf{Detection:}
   Global forecast data from the German Weather Service in NetCDF format is processed by the existing CGNet-based segmentation network. For each of the timesteps, the network produces probability fields for background, tropical cyclone (TC), and atmospheric river (AR). This detection component was developed as prior work and is not part of this study.

    \item \textbf{Feature Extraction and Tracking:}
    Each timestep's probability field is thresholded and segmented independently into candidate features, which are then linked across consecutive timesteps into continuous tracks.

    \item \textbf{Feature Ranking and Motion-Type Selection:}
   Each track is characterized using four measures and assigned a relevance score. Based on these characteristics, one of three camera-motion types is selected: follow, skeleton, or orbit.

    \item \textbf{Camera Trajectory Generation and Rendering}
    The ranked, motion-typed tracks produced by stage 3 are converted into Met3D-compatible camera keyframe sequences (Section 2.3), which are then rendered and validated against Met3D's visual output (Section 3).
    
\end{enumerate}

\section{Evolution of the Design}

The final design (Section 5) was developed through four iterations. Each
iteration exposed a limitation that shaped the next one.

\subsection{Unselective Coverage and Point-Hopping}

An early design treated every detected TC and AR as a point of interest.
The camera moved through these points in the order they were detected,
without any selection or geographic reasoning. A single forecast can
contain dozens of active features at once. This produced an animation of
roughly 1000 frames, in which the camera jumped between widely separated
points across the globe.

Fig.~\ref{fig:point-hopping} shows this behaviour at a single forecast
timestep. The camera path crosses ocean basins and continents with no
coherent structure. At one point, it jumps directly from the North
Atlantic to Southeast Asia between two consecutive waypoints. The large
number of transitions, combined with how abrupt they were, made the
motion difficult to follow. When several features were active at once,
the viewer could not tell which one was currently being shown. Adding
more frames did not make individual features easier to follow.

This iteration showed that displaying every detected feature is not the
same as displaying something meaningful. The design first needed to
decide which features were worth showing, before deciding how to move
between them.

\begin{figure}[t]
    \centering
    \includegraphics[width=\linewidth]{figures/camera_path_00.png}
    \caption{Camera path at a single forecast timestep in the
    point-hopping design, visiting numbered waypoints attached to
    individually detected TC (red) and AR (teal) features in the order
    they were found. The path spans large geographic distances between
    consecutive waypoints. For instance, it jumps directly from the North
    Atlantic (waypoint 28) to Southeast Asia (waypoint 25), rather than
    following a coherent geographic or narrative order.}
    \label{fig:point-hopping}
\end{figure}

\subsection{Constrained and Manually Authored Alternatives}

Two more constrained alternatives were tested next. Neither yet ranked
features by importance.

The first restricted the camera to a fixed region, the North Atlantic
corridor. The assumption was that narrowing the domain would make the
task easier. This assumption did not hold. Relevant TCs and ARs often
formed or moved outside the region. At times, the region itself was
inactive while significant activity occurred elsewhere, leaving the
camera with nothing to show.

The second alternative covered features in a fixed geographic order,
from north to south, using manually authored Met3D XML sequences. This
avoided the coverage problem of the first alternative. However, it
required considerable manual effort for every new forecast. It also
produced inconsistent camera parameters between features. Beyond about
five features, the approach became impractical.

Alongside this manual scripting, a fully hand-authored reference
animation was produced early in the project. This animation served two
purposes. It established a target look for the eventual system, and it
helped identify Met3D's pitch and yaw conventions, which are only
lightly documented in the application itself. This exploratory work
later informed the camera-calibration procedure described in
Section 5.3.5.

Together, these two alternatives showed that neither a fixed region nor
a manually authored sequence could generalise across an arbitrary
forecast. This motivated a shift toward automated, data-driven feature
ranking. The regional-restriction idea was not discarded completely. It
was later kept as an optional camera grammar, a regional cut applied on
top of the final ranking rather than a fixed default (Section 5.3.4).
This allows corridor-based framing when needed, without limiting the
coverage of the full system.

\subsection{Weighted Multi-Criteria Scoring}

The third design abandoned manual scripting and fixed regions.
Instead, it automatically selected a small number of top-priority
tracks, and only then decided how to film them.

Each candidate track was scored using a weighted combination of four
attributes. Average speed (35\%) prioritised fast-moving systems that
would produce dynamic camera motion. System lifetime (25\%) favoured
long-lasting, stable phenomena over transient noise. Split/merge
ambiguity (25\%) highlighted complex interaction points where
neighbouring features diverge or converge. Area growth rate (15\%)
identified rapidly intensifying systems.

This approach solved the coverage problems of the earlier designs. It
introduced a new one, however. The composite score showed that a track
ranked highly, but not why it ranked highly. A top-ranked track could
be fast, long-lived, volatile, or simply average across all four
attributes at once, and the score alone could not distinguish between
these cases. Each camera-motion type depends on a feature having one
clearly dominant trait. This made it impossible, in general, to decide
how a top-ranked track should be filmed.

\subsection{Dominant-Metric Prioritization and Cinematic Rendering}

The final design solves this problem directly. Instead of a weighted
composite, it compares three metrics separately: travel distance,
spatial extent, and shape volatility. Each metric is normalised
independently. A track's score and its camera-motion type are both
determined by whichever metric it is most exceptional in, rather than
by a blended total. This guarantees that every selected track has one
identifiable, dominant characteristic to build a camera motion around
(Section 5.2).

The selected features are then rendered using a small set of cinematic
techniques. These include zoom, pitch adjustment, and a pause-and-spin
motion, in which the forecast is held still while the camera circles a
feature. This technique was inspired by NASA's atmospheric-river
visualisations \citep{nasa-climate}, which similarly pause time to
highlight a feature of interest.

Compared with the two preceding designs, the resulting animations were
substantially more engaging. Viewers could follow how a feature evolved
and changed shape, without being distracted by unrelated activity
elsewhere in the forecast. 

\section{Methodology}

\subsection{Feature Extraction and Tracking}

\textbf{Input representation:} The method operates on a sequence of NetCDF probability fields $P_t(x,y)$, one per forecast timestep $t \in \{1,2,3,\dots\}$, each giving the network's estimated likelihood that location $(x,y)$ belongs to the feature class of interest (TC or AR).

\textbf{Thresholding:} Each probability field is converted to a binary candidate mask by thresholding at $\tau = 0.5$:
\begin{equation}
B_t(x,y) =
\begin{cases}
1, & P_t(x,y) \geq \tau \\
0, & P_t(x,y) < \tau
\end{cases}
\tag{1}
\end{equation}

\textbf{Connected-component filtering:} Connected-component labelling is applied to $B_t$ to identify spatially contiguous candidate features. Components smaller than a class-specific minimum size are discarded to remove noise and isolated detections. The current implementation uses minimum sizes of 50 pixels for TC candidates and 200 pixels for AR candidates, reflecting the different spatial scales of the two feature types.

\textbf{Spatial descriptors:} For each surviving component, a set of geographic descriptors is computed. For tropical cyclones, the descriptor of interest is the centroid, taken as the center of mass of the component's pixels:
\begin{equation}
\mathbf{c} = \frac{1}{n}\sum_{j=1}^{n}(x_j,y_j),
\qquad
(x_j,y_j) \in \Omega
\tag{2}
\end{equation}
where $\Omega$ is the set of pixels belonging to the component.

For atmospheric rivers, no single point is representative, and the component is instead reduced to its skeleton which is a one-pixel-wide centerline obtained via a standard skeletonization algorithm. The skeleton's two endpoints and its midpoint serve as the key waypoints characterizing the river's spatial extent, analogous to finding the backbone of a ribbon-shaped object.

\textbf{Temporal tracking:} Each timestep's set of feature observations $F_t = \{f_1,\dots,f_n\}$ is associated with the observations at $t+1$ using nearest-neighbour matching on the great-circle (Haversine) distance between centroids. For two points at latitude/longitude $(\varphi_1,\lambda_1)$ and $(\varphi_2,\lambda_2)$, expressed in radians:
\begin{equation}
a =
\sin^2\!\left(\frac{\Delta\varphi}{2}\right)
+
\cos\varphi_1\cos\varphi_2
\sin^2\!\left(\frac{\Delta\lambda}{2}\right),
\qquad
d =
2R \cdot
\operatorname{atan2}\!\left(\sqrt{a},\sqrt{1-a}\right)
\tag{3}
\end{equation}
where $R \approx 6371\,\text{km}$ is the Earth's mean radius.A feature at time $t$ is linked to the closest feature at $t+1$ when the distance $d$ is within a predefined maximum radius. This radius is used to allow for the expected movement of a feature between two consecutive forecast hours. If the distance is larger than this radius, the two features are not linked, which helps avoid connecting unrelated detections. 

The resulting associations form a trajectory for each feature $i$:
\begin{equation}
T_i = \{f_{i,t_1}, f_{i,t_2}, \dots, f_{i,t_k}\}
\tag{4}
\end{equation} represents the life cycle of the feature over the timesteps in which it was detected.

\subsection{Feature Ranking}

\textbf{Characterization:} Each completed trajectory $T_i$ is summarized along four metrics intended to capture how relevant the feature is likely to be for visualization.

\textit{Lifetime} is the duration over which the feature could be continuously tracked:
\begin{equation}
L_i = t_k - t_1
\tag{5}
\end{equation}

\textit{Total travel distance} is the cumulative geographic displacement of the feature's centroid over its tracked lifetime, computed as the sum of consecutive Haversine distances (Eq. 3):
\begin{equation}
D_i =
\sum_{j=1}^{k-1}
d\!\left(
\mathbf{c}_{i,t_j},
\mathbf{c}_{i,t_{j+1}}
\right)
\tag{6}
\end{equation}

\textit{Spatial extent} $E_i$ describes the geographic size of the feature, derived from the spatial dimensions (e.g., bounding-box diagonal or mask area) of the detected component, averaged or maximized over the track.

\textbf{Shape volatility.} $V_i$ combines two signals: how much the
feature's area changes frame-to-frame, and how often it has a nearby
same-class neighbour. In plain terms: for every timestep the track was
observed, the procedure checks whether any other detected blob of the same
feature class lies within a radius $r_{sm}$; if one does, that timestep is
counted as one split/merge interaction. The two signals are the mean
frame-to-frame change in pixel area,
\begin{equation}
\overline{\Delta A}_i =
\frac{1}{k-1}
\sum_{j=1}^{k-1}
\left| a_{i,t_{j+1}} - a_{i,t_j} \right|
\tag{8a}
\end{equation}
and the count of timesteps with at least one nearby same-class neighbour,
\begin{equation}
n_i^{\mathrm{sm}} = \sum_{j=1}^{k} \delta_j,
\qquad
\delta_j =
\begin{cases}
1, & \text{a same-class blob lies within } r_{sm}\text{ km of } f_{i,t_j} \\
0, & \text{otherwise}
\end{cases}
\tag{8b}
\end{equation}
where $\delta_j$ is simply a per-timestep flag: one interaction, one count,
nothing more elaborate. The two signals are combined into a single composite
score using a weighting constant $w_{sm}$ that converts each interaction
event into area-change-equivalent units:
\begin{equation}
V_i = \overline{\Delta A}_i + w_{sm}\, n_i^{\mathrm{sm}},
\qquad w_{sm} = 500
\tag{8}
\end{equation}
\textbf{Normalization and scoring:} Because the four metrics have incompatible units and scales, each is standardized to a $z$-score across all tracked features in the dataset before combination:
\begin{equation}
Z_{X,i} =
\frac{X_i-\mu_X}{\sigma_X},
\qquad
X \in \{D,E,V\}
\tag{8}
\end{equation}
where $\mu_X$ and $\sigma_X$ are the mean and standard deviation of metric $X$ over all tracks. The normalized metrics are combined into a single visualization-relevance score via a weighted sum:
\begin{equation}
S_i =
w_D Z_{D,i}
+
w_E Z_{E,i}
+
w_V Z_{V,i}
\tag{9}
\end{equation}

Tracks are ranked by $S_i$, and the highest-scoring features are selected for camera animation. The weights $w_L,w_D,w_E,w_V$ are treated as a visualization-oriented design choice rather than an objective measure of meteorological importance. Different weightings emphasize different kinds of ``interesting'' behaviour and their sensitivity is intended to be examined empirically.

\textbf{Camera-motion-type selection:} The same characterization that drives ranking also determines \textit{how} a selected feature should be presented. Three camera-motion archetypes are defined:

\begin{table}[htbp]
\centering
\caption{Mapping from dominant feature characteristic to camera-motion archetype.}
\label{tab:camera-motion}
\begin{tabular}{lll}
\hline
\textbf{Dominant characteristic} &
\textbf{Camera-motion type} &
\textbf{Visualization objective} \\
\hline
High travel distance &
Follow &
Emphasize the feature's geographical movement \\

Large spatial extent &
Skeleton &
Reveal the feature's spatial scale and structure \\

High shape volatility &
Orbit &
Reveal changes in the feature's morphology \\
\hline
\end{tabular}
\end{table}

This mapping ties the camera's behaviour to a measurable property of the tracked phenomenon, rather than choosing a motion independently of the data: a fast-translating cyclone earns a chase-style follow shot; a continent-spanning atmospheric river earns a shot that traces its skeleton to convey scale; a rapidly reorganizing storm earns an orbit that lets the viewer see structural change from multiple angles.

\textbf{Output interface:} For each selected feature, its trajectory (Eq. 4), the four characterization metrics (Eqs. 5--8), the composite score (Eq. 9), and the assigned motion type are serialized to an intermediate JSON file, along with supporting spatial information such as bounding boxes, needed by the camera-generation stage. This file is the sole interface to Section 4.3.


\subsection{Camera Path Generation}

Given the ranked, motion-typed feature list from Section 2.2, this stage
generates the Met3D camera keyframe sequences. Each ranked feature is
assigned one of three motion types: skeleton, follow, or orbit
(Table~\ref{tab:camera-motion}). The following subsections describe how
each motion type is implemented as a concrete camera behaviour.

\subsubsection{Skeleton Animation}

For features assigned the skeleton motion type, the camera path follows
the feature's morphological skeleton rather than its centroid. The path
is anchored at the timestep where the feature reaches its maximum
spatial extent. At this timestep, the feature's skeleton (Section 2.1)
is re-extracted to determine its primary axis of elongation. This
skeleton then defines the waypoints for the full camera sequence.

Skeletonization was chosen on the recommendation of the Met3D developers
(C. Fischer, personal communication), since it reliably captures the
morphological spine of irregular AR shapes. Unlike a centroid-based
path, the skeleton preserves the feature's elongation. This allows the
camera to travel along the length of the feature, rather than toward a
single point.

To avoid the camera clipping into the feature or producing blind spots,
a fixed spatial offset is applied between each skeleton waypoint and
the corresponding camera position:
\begin{equation}
\varphi_{\mathrm{cam}} = \varphi_{\mathrm{AR}} - \Delta\varphi_{\mathrm{off}},
\qquad
\lambda_{\mathrm{cam}} = \lambda_{\mathrm{AR}} - \Delta\lambda_{\mathrm{off}}
\tag{10}
\end{equation}
$\Delta\varphi_{\mathrm{off}}$ and $\Delta\lambda_{\mathrm{off}}$ are set
to $5^{\circ}$ south and $5^{\circ}$ west of the feature. This value was
found through iterative testing in Met3D. The offset moves the camera
away from the feature, while the yaw computation below keeps the camera
pointed toward it.

At each waypoint, the camera yaw is computed so that the camera points
toward the next waypoint along the skeleton path:
\begin{equation}
\psi_t = -\operatorname{atan2}\!\left(\lambda_{t+1}-\lambda_t,\ \varphi_{t+1}-\varphi_t\right)
\tag{11}
\end{equation}
The negative sign keeps the camera facing along the west-to-east
direction of the feature's elongation. This direction matches the
predominant orientation of the AR skeletons observed in the evaluation
data.

As the camera moves along the skeleton, its pitch is varied to show the
feature's vertical extent and structure. Pitch decreases from
$60^{\circ}$ at the start of the traversal to $20^{\circ}$ at the end.
This produces a pull-back effect as the camera leaves the feature.

\subsubsection{Follow Animation}

\textit{Pending. The material available so far describes only the
skeleton motion type in detail. The implementation of the follow
motion type, used for features whose dominant characteristic is
travel distance, is not yet documented here.}

\subsubsection{Orbit Animation}

\textit{Pending. The material available so far describes only the
skeleton motion type in detail. The implementation of the orbit
motion type, used for features whose dominant characteristic is
shape volatility, is not yet documented here.}

\subsubsection{Algorithm Summary}

The NetCDF probability fields and the ranked-feature JSON (Section 2.2)
are combined to locate the timestep of maximum spatial extent for each
selected feature. The feature's skeleton is then extracted at that
timestep. Waypoints are selected along the skeleton, or along the
centroid track for follow features. The resulting waypoint sequence,
together with the offset, yaw, and pitch parameters described above, is
written to a Met3D-compatible XML keyframe file.
\section{Results and Evaluation}

\subsection{Evaluation Data}

Three datasets were used to develop and test the full system, from
feature extraction through to camera-path generation. The first is a
12-timestep forecast dataset, described in Section 5.1. The second,
named Dublin, spans 20 to 27 January, at 12:00 UTC. The third, named
Limerick, comprises 61 timesteps over 180 hours (7.5 days), from 28
January to 4 February 2026.

\subsection{Validation Methodology}

Validation was carried out through iterative manual inspection in Met.3D. The process consisted of generating a candidate XML file from the Python script, loading it into Met.3D, and playing the resulting animation. The camera behaviour was then inspected to determine whether the camera followed the skeleton path correctly, whether the zoom, pitch, and yaw values were appropriate for the scene, and whether the time advancement matched the evolution of the data.

If the behaviour was not satisfactory, the parameters were manually adjusted using Met.3D's camera editor. The corrected XML was then exported and used to update the corresponding values in the Python code. The script was subsequently re-run and the resulting animation was validated again.

This iterative process was used to check visual correctness and to identify parameter ranges that produced suitable animations. The qualitative validation criteria were that the camera followed the skeleton path without major deviations, that the zoom provided good feature visibility without clipping, that the pitch and yaw produced a natural cinematic effect, and that the time evolution showed feature growth or dissipation as expected.

\subsection{Comparative Analysis: Manual versus Automated Generation}

\begin{table}[htbp]
\centering
\caption{Comparison of manual and automated camera-sequence generation.}
\label{tab:manual-vs-automated}
\begin{tabular}{lccc}
\hline
\textbf{Metric} & \textbf{Manual approach} & \textbf{Automated approach} & \textbf{Improvement} \\
\hline
Time per AR & 20--60 min & 2--5 s & $\sim$95\% reduction \\
Parameter consistency & High variability & Standardized defaults & High uniformity \\
Scalability & Limited (2--3 ARs) & Extensive (10+ ARs) & 5--10$\times$ scaling \\
Reproducibility & Manual, error-prone & Deterministic, code-based & Fully reproducible \\
Customization & Ad-hoc & Systematic (JSON input) & Flexible, data-driven \\
\hline
\end{tabular}
\end{table}

The automated approach reduces the time required to generate a single AR camera sequence by roughly two orders of magnitude compared with manual authoring. It also provides more consistent parameter choices and allows the generation process to scale to substantially more simultaneous features than is practical by hand.

\subsection{Challenges Encountered and Resolution}

\textbf{Parameter interdependence.}
Zoom, pitch, and yaw are not independent, and a poor choice of one parameter can affect the usefulness of the others. For example, a steep pitch of $60^{\circ}$ at moderate zoom (level 100) can cause clipping, while a shallow pitch of $20^{\circ}$ at close zoom (level 25) does not convey the vertical structure well. Through approximately 20 iterative test cycles in Met3D, workable parameter ranges were identified. Close zoom (level 25) requires a high pitch of $60^{\circ}$ to avoid ground clipping. Regional zoom (level 100) works best with a moderate pitch of $30$--$40^{\circ}$, while global zoom (level 250) requires a low pitch of $0$--$10^{\circ}$ to maintain geographic context. These ranges were encoded as defaults in the code, with optional overrides through the input JSON.

\textbf{Skeleton quality variability.}
Features with irregular or fragmented shapes can produce poor skeletons, including disconnected segments or spurious branches. These can result in unrealistic camera paths. This was addressed using the minimum component-size threshold introduced in Section 2.1 (50 pixels) and by interpolating waypoints across small gaps in the skeletonized region.

\textbf{Timeline non-reversibility.}
Met3D's time axis advances monotonically and cannot rewind to $t=0$ after reaching the final timestep. This prevents the full life cycles of multiple features from being shown within a single continuous animation. To address this, separate XML files are generated for each feature when the combined lifetime exceeds the available number of frames, which is 61 for the 180-hour evaluation dataset. The resulting files are loaded sequentially by the user.

\textbf{Splitting and merging.}
Some features split into multiple disconnected blobs or merge with neighbouring systems. The current skeleton-extraction step assumes a single continuous object and therefore produces multiple fragmented paths in these cases. This issue is currently unresolved. Split and merge events are recorded in the feature metadata (Eq. 7b), but a corresponding camera-visualization strategy has not yet been developed. Multi-object tracking or temporal skeleton interpolation are identified as possible directions for future work (Section 5).


\subsection{Limitations}

The current implementation has several limitations:
\begin{enumerate}
\item It assumes a single continuous feature object and cannot yet handle splitting or merging events during camera-path generation.
\item Manual validation of the generated XML is still required, and no automated visual-quality score is currently available.
\item Some pitch and zoom values are hardcoded, which limits user customisation beyond the JSON override mechanism. \cite{caban2007texture}
\item Tracking correctly follows features across the ±$180^{\circ}$ longitude line, but the camera path itself does not account for this wraparound, causing it to jump abruptly from one edge of the map to the other.
\end{enumerate}


% %% REFERENCES

% %% The reference list is compiled as follows:
\bibliographystyle{copernicus}
\bibliography{bib}

\end{document}
